\input amstex
\documentstyle{tran}
\title
On Trace Identities and Universal Eigenvalue Estimates 
for Some Partial Differential Operators\endtitle

\author Evans M. Harrell II and Joachim Stubbe\endauthor

\address School of Mathematics, 
Georgia Institute of Technology, 
Atlanta  GA  30332-0160\endaddress

\email harrell\@math.gatech.edu\endemail


\address D\'epartement de Physique Th\'eorique,
Universit\'e de Gen\`eve,
Geneva, Switzerland\endaddress     

\email stubbe\@cernvm.cern.ch\endemail

\cvol{00}
\cvolno{00}
\cvolyear{0000}
\cyear{0000}
\cmonth{XXX}

\subjclass 35J10, 35J25, 58G25\endsubjclass  

%%version of 17 September 1995
\abstract
In this article, we prove and exploit a trace identity for the spectra of 
Schr\"odinger operators and similar operators.  This identity leads to 
universal bounds on the spectra, which apply to low-lying eigenvalues, 
eigenvalue asymptotics, and to partition functions (traces of heat operators).  
In many cases they are sharp in the sense that there are specific 
examples for which the inequalities are saturated.  
Special cases corresponding to known 
inequalities include those of Hile and Protter and of Yang.
\endabstract

\thanks The first author
was supported in part by US NSF Grant DMS 9211624.\endthanks

\date {} \enddate
\keywords Schr\"odinger operator, eigenvalue gap, trace,
heat kernel, partition function\endkeywords
\endtopmatter

\document
\head Introduction\endhead

In this article, we prove and exploit an identity for the 
spectra of self-adjoint 
operators $H$ modeled on the Dirichlet Laplacian or, more generally, on 
Schr\"odinger operators of the form
$$
(\bold p - \bold A(\bold x))^2 + V(\bold x),                                    
\tag 1$$
where ${\bold p} = {1\over i}\, \nabla$ 
is the usual momentum operator in convenient 
units and {\bf A(x)} is the magnetic vector potential.   
We recover and extend 
several known inequalities involving sums, differences, and ratios of 
eigenvalues.  Let 
$\lambda_{j}$, $j = 1,\dots,$ denote the ordered eigenvalues 
of the Dirichlet Laplacian on a bounded 
$d$-dimensional domain with zero 
Dirichlet boundary conditions, and recall that Hile and Protter [HiPr80] 
proved that:
$${d \over 4} \le  {1 \over n}\sum\limits_{j = 1}^{n} 
{{\lambda }_{j} \over {\lambda }_{n+1} - {\lambda 
}_{j}},\tag 2 $$
thereby extending an earlier inequality of Payne, P\'olya, and 
Weinberger\break 
[PaPoWe56].  In the last few years it has become clear that these and many 
similar relationships can be realized as special 
cases of abstract variational 
bounds involving the interplay among commutators of $-\nabla^2$, a 
Cartesian coordinate $x_j$ and the corresponding derivative 
$\partial/\partial x_j$.  
For this analysis and various extensions of (2) see [Ha88], 
[Ho90], [Ha93], [HaMi95].  

While inequalities of this type have been fairly sharp for low-lying 
eigenvalues, they have been mostly disappointing for higher eigenvalues.  
Yang [Ya91], however, first noticed that it is possible to use a modification 
of the original Hile-Protter argument to obtain somewhat complicated 
inequalities, the two sides of which are of the same order for high 
eigenvalues, when compared with the Weyl asymptotic formula for the 
Dirichlet Laplacian.  In the course of understanding Yang's results, we 
discovered an exact identity which implies Yang's bound, the Hile-Protter 
bound (2), and several new ones.  This identity is the cornerstone of our 
analysis.  We emphasize that it is entirely nonvariational.

\head I.  An abstract trace formula\endhead

The operators to which our identity applies are modeled on (1) but 
potentially more general.  Let $H$ be a self-adjoint operator with purely 
discrete spectrum on a Hilbert space with inner product denoted 
$\langle,\rangle$.  Its 
eigenvalues and its normalized eigenfunctions will be written as 
$\lambda_{j}$ and $u_j$, respectively, $j = 1,2,\dots$.  There will be a 
distinguished subset $J$ of the spectrum of $H$; $J^c$ 
will denote its complement 
in the spectrum, and $P_J$ and $P_{J^c}$ will be the corresponding spectral 
projections.  We shall be interested in traces of  
$P_J f(H)$, where $f(z)$ is any 
function defined on the spectrum of $H$.

\proclaim{Theorem 1}  
Suppose there is a family of symmetric operators $X_{\alpha}$ 
and skew symmetric operators  $D_{\alpha}$, 
$\alpha = 1, \dots, d$, such that  
$[D_{\alpha},X_{\alpha}] = 1$,  and that $H$ is of the form
$$
H =   - \sum_{\alpha=1}^d D_{\alpha}D_{\alpha}    +      
V,   \tag 3
$$
where $V$ is such that  $[V,X_{\alpha}] = 0$.  
Furthermore, suppose that 
$X_{\alpha}u_j$ is in the domain of $H$ for any 
$j$ and that $D_{\alpha}$ is 
$H$-bounded.  Define 
$$T_{\alpha jm} := |(D_{\alpha} u_j, u_m)|^2    \ge   0.$$
Then:
$$
\sum\limits_{\lambda_j\in J} f\left(\lambda_j\right)
= -2\sum\limits_{\scriptstyle\lambda_j,\lambda 
_m\in J\atop
\scriptstyle\lambda_j\ne \lambda_m}
{f\left(\lambda _j\right) - f\left(\lambda_m\right) 
\over\lambda_j-\lambda _m}\,
{T}_{\alpha jm}+4\sum\limits_{\scriptstyle\lambda _j\in J\atop
\scriptstyle\lambda_q\in J^c}
{f\left({{\lambda }_{j}}\right) 
\over {\lambda }_{q}-{\lambda }_{j}}\,{T}_{\alpha jq} , 
\tag 4 $$
provided that these sums converge.
\endproclaim 

\remark{Remarks}  %% if EH wants plural use demo{}
\roster
\item Multiple eigenvalues are counted with their degeneracies.
\item In our principal model (1), 
$D_{\alpha} = (\nabla - i {\bold A})_{\alpha}$, 
and 
$$T_j := \Sigma _{\alpha m} T_{\alpha m j}                                                                                     
\tag 5$$
corresponds to the kinetic energy of the state $u_j$.  
Moreover, in applications 
where 
$H = - \nabla ^2$, $T_j$ 
is the usual Dirichlet integral with $u_k$, so $T_j = 
\lambda_{j}$.  The same is true for purely magnetic Schr\"odinger 
operators, $H = -(\nabla - i {\bold A})^2$.  
\item There are other straightforward extensions, 
which we shall not bother to 
write explicitly, such as replacing the Laplacian with other constant 
coefficient elliptic operators.  
Likewise, we note that the formula extends to 
the case of continuous spectra when the density of 
states is well defined.  In 
the latter case the sums become integrals, which must be assumed 
convergent. 
\item If $H$ is an abstract operator not necessarily satisfying (3), and 
$[D_{\alpha} ,X_{\alpha}] = 1$ for a particular pair of operators 
$D_{\alpha}$, 
$X_{\alpha}$, we still obtain the identity (6) given below.
\item This identity was suggested by the analysis in [Ya91], 
where there is an 
inequality related to the case 
$J = \{\lambda_1, ..., \lambda_n\}$, 
$f(\lambda_j) = (\lambda_{n+1} - \lambda_j)^2$, $H = -\nabla^2$.  
(See Theorem 5, below.)  
\endroster
\endremark

\demo{Proof}
We rewrite tr$(P_J f(H))$ as follows:
$$\eqalign{
\sum\limits_{{\lambda }_{j}\in J} 
f(\lambda _j) &= \text{tr}
(P_J f(H))=\text{Tr}
(P_Jf(H)[D_\alpha, X_\alpha ])\cr
&=\text{tr}(P_J f(H) [D_{\alpha} (P_J + P_{J^c}) X_{\alpha} 
-X_{\alpha} (P_J + P_{J^c}) D_{\alpha}] )\cr
&=\text{tr}(P_J f(H) [D_{\alpha} P_J X_{\alpha} - X_{\alpha} P_J 
D_{\alpha}] )\cr   
&\qquad+ \text{tr}(P_J f(H) [D_{\alpha} P_{J^c} X_{\alpha} - X_{\alpha} P_{J^c} 
D_{\alpha}] ).
}$$
Using the gap formula 
$$(\lambda_{m} - \lambda_{j}) \langle X_{\alpha}u_j,u_m\rangle
=\langle [H,X_{\alpha}]u_j,u_m\rangle 
$$
we obtain
$$\eqalign{
\text{Tr}(P_Jf(H)&(D_{\alpha}P_JX_{\alpha} - X_{\alpha}P_JD_{\alpha}))
= -2\sum\limits_{\scriptstyle\lambda_j,\lambda_m\in J\atop
\scriptstyle\lambda _j\ne \lambda_m}
f(\lambda_j)\langle X_{\alpha}u_j, 
u_m\rangle\langle u_m,D_{\alpha }u_j
\rangle\cr
&= 2 \sum\limits_{\scriptstyle\lambda _j, 
\lambda_m\in J\atop
\scriptstyle
\lambda_j\ne \lambda _m}{f(\lambda_k)\langle[H,X_{\alpha }]u_j , 
u_m\rangle\langle u_m,D_{\alpha }u_j 
\rangle \over \lambda _j -\lambda_m}\cr
&=\sum\limits_{\scriptstyle\lambda_j ,\lambda 
_m\in J\atop
\scriptstyle\lambda_j\ne\lambda _m}{f(\lambda 
_j)- f(\lambda _m) \over \lambda _j - 
\lambda_m}\langle [H,X_\alpha ]
u_j,u_m\rangle\langle u_m, 
D_\alpha u_j\rangle,}
$$
where in the last step we symmetrized in the two indices.  
The cross term 
connecting $J$ and $J^c$ 
can be rewritten similarly, except that symmetrization 
does not apply.  The result is 
$$\eqalign{\sum\limits_{\lambda_j\in  J} 
f(\lambda_j)&=
 \sum\limits_{\scriptstyle\lambda _j ,\lambda_m\in 
J\atop
\scriptstyle\lambda_j\ne \lambda _m}{f(\lambda 
_j) - f(\lambda _m) \over\lambda_j- 
\lambda_m}\,\langle[H,X_\alpha]
u_j,u_m\rangle\langle u_m,D_\alpha u_j\rangle\cr
&\qquad +\sum\limits_{\scriptstyle\lambda_j\in J\atop
\scriptstyle\lambda_q\in J^c}{2f(\lambda _j 
) \over \lambda_j-\lambda_q}\,\langle[
H,X_\alpha]u_j,u_q\rangle
\langle u_q,D_\alpha u_j
\rangle.} \tag 6$$
Identity (4) then follows because $[H,X_{\alpha}] = - 2 D_{\alpha}$.
\hfill\qed
\enddemo

With the choice $f(\lambda) = (z - \lambda)^2$, identity (4) can be 
simplified, as follows:

\proclaim{Proposition 2}
With the same assumptions as in \rom{Theorem 1}, 
$$
d\sum\limits_{\lambda _j\in J}({z 
-\lambda_j})^{2} =4\sum\limits_{\lambda_j\in J}
(z -\lambda_j)T_j+4\sum\limits_{\scriptstyle\lambda_j\in J\atop
\scriptstyle\lambda_q\in J^c}{(z -\lambda_j)
(z-\lambda_q) \overwithdelims () \lambda_q-\lambda_j}
\sum\limits_\alpha T_{\alpha 
jq}.\tag 7$$
\endproclaim

\demo{Proof}  
The first term on the right in (4) is
$$\eqalign{
-2\sum_{\scriptstyle\lambda_j,\lambda_m\in J\atop 
\scriptstyle\lambda_j\ne\lambda_m}
{(z-\lambda_j)^2-(z-\lambda_m)^2\over \lambda_j-\lambda_m}
\,T_{\alpha jm} &= 2\sum_{\lambda_j,\lambda_m\in J}
(2z-\lambda_j-\lambda_m)T_{\alpha jm}\cr
&=4\sum_{\lambda_j,\lambda_m\in J}(z-\lambda_j)T_{\alpha jm}}$$
(because of the symmetry in $j,m$)
$$ = 4\sum\limits_{\lambda_j\in J}(
z -\lambda_j)\sum\limits_{\hbox{all }\lambda_m} 
T_{\alpha jm}-4\sum\limits_{\scriptstyle\lambda_
j\in J\atop\scriptstyle\lambda_m\in J^c}
(z -\lambda_j)T_{\alpha jm}.$$
Equation (7) now results from summing (4) over 
$\alpha$, using (5), and noting 
that
$$
{(z-\lambda_j)^2 \over \lambda_q- 
\lambda_j}- (z-\lambda_j) = {(z
-\lambda_j)(z -\lambda_q) \over \lambda 
_q -\lambda_j}.\eqno{\qed} $$
\enddemo

\head II.  The estimates of Hile-Protter and Yang, and extensions\endhead


Earlier generalizations of the Hile-Protter inequality (2) in 
[Ha88], [Ho90],\break 
[Ha93], [HaMi95] had no analogue of the function $f$, 
and were on the other 
hand based on variational estimates of eigenvalues.  
The proof of Theorem 1 
is not only more direct than the earlier estimates,  
but it also allows several 
new corollaries.  
In addition, several of our corollaries are sharp; 
this will be 
indicated by the use of the symbol \#.

\proclaim{Corollary 3} (\#, harmonic oscillator, $n=1$, all $d$).  
$$
{d \over 4}\le{1 \over n}\sum\limits_{j \le n} 
{T_j \over \lambda_{n+1} -\lambda_j}.\tag 8$$
\endproclaim

\remark{Remarks}
This includes Hile-Protter (2), since for the Dirichlet Laplacian 
$T_j = \lambda_{j}$.  Because of the virial theorem, for Schr\"odinger 
operators there are often substantially better substitutes 
for $T_j$  than the 
inequality $T_j   \le  \lambda_j$, and in this case (8) is sharp 
(see Corollary 3).  
The Hile-Protter inequality for the Dirichlet Laplacian is 
definitely not sharp, even for the lowest two eigenvalues, for which the 
sharp version is that of Ashbaugh and Benguria [AsBe].  A generalization of 
inequality (8) is given below in Theorem 5.
\endremark

\demo{Proof}
Make the choice $f(z) = 1$, $J = \{\lambda_1, ..., \lambda_n\}$, 
in (4) and 
sum over all $\alpha$ from 1 to $d$.  
The first term on the right vanishes and the sum 
in the second term can be extended to all $q$ to yield:
$$
 d n = 4\sum\limits_{{\scriptstyle\alpha \atop
\scriptstyle j\le n}\atop
\scriptstyle q \ge n+1} {T_{\alpha jq} \over \lambda_
q-\lambda_j}\le 4 \sum\limits_{{\scriptstyle\alpha \atop
\scriptstyle j \le n}\atop\scriptstyle q \ge 1}{T_{\alpha jq} \over \lambda_
{n+1} - \lambda_j} = 4 \sum\limits_{j\le n} 
{T_j \over \lambda_{n+1} -\lambda_j}. 
 $$
\enddemo
	
In at least one important situation, inequality (8) is sharp.  
For consider the 
quantum harmonic oscillator in $d$ dimensions, 
$$
H = |{\bold p}|^2 + |{\bold x}|^2,
$$
the lowest two eigenvalues of which are known to be 
$\lambda_1 = d$ and 
$\lambda_2 = d+2$. Because of the virial theorem, $T_j = {1\over 2} 
\lambda_{j}$ for the harmonic oscillator, and therefore with $n=1$ the right 
side of (8) equals  ${d\over 2}{1\over 2} ={d\over4}$, and the 
inequality is saturated.  The 
general situation applying the virial theorem is: 

\proclaim{Corollary 4} (\#, harmonic oscillator, $n=1$, all $d$).  
Let $H =   |{\bold p}|^2 + V({\bold x})$ on 
$L^2(R^d)$, where we assume that
$${\bold x}\cdot \nabla V   \le  \beta V,          
\tag 9$$
and suppose that $\lambda_1, \dots \lambda_{n+1}$ 
are discrete eigenvalues.  
Then 
$$
{d(2+\beta) \over 4\beta }\le{1 \over 
n}\sum\limits_{k\le n}{\lambda _k \over \lambda 
_{n+1} -\lambda_k}.$$
\endproclaim

\remark{Remarks}
This applies in particular to power-law potentials $V({\bold x}) =  
|{\bold x}|^{\beta}$, $\beta > 0$, 
and even, despite the presence of continuous 
spectrum, for the negative eigenvalues with 
$V({\bold x}) = - |{\bold x}|^{\beta}$, 
$-2 < \beta < 0$.  The correct generalization 
corresponding to $\beta =0$ is 
$V({\bold x}) = \ln(|{\bold x}|)$, for which the conclusion reads
$${d \over 2} \le  {1 \over n}\sum\limits_{k\le n}{
1 \over \lambda_{n+1} - \lambda_k}. $$
Obviously, any of these power-law potentials could be multiplied by a 
positive constant with a similar estimate.
\endremark

\demo{Proof}  
By the virial theorem (see [EaKa]), 
$$T_j=
{1\over 2} \langle u_j,({\bold x}\cdot\nabla V)u_j\rangle\le{\beta\over
2}\,\langle u_j,Vu_j\rangle.$$
Since  ${\lambda }_{j} ={T}_{j} +\langle u_j, V 
u_j\rangle,$ a simple calculation shows that this implies that 
$${T}_{j}\le \sigma {\lambda }_{j },
	 \tag 10$$
with $\sigma=\beta /(2+\beta )$, which when inserted into (8) yields 
the claim. 
\hfill\qed\enddemo


We next present some inequalities related to 
identity (4) for other choices of 
the function $f$.  
The important new idea in Yang's analysis, as we interpret 
it, was to change the $f$ 
used by Hile and Protter, so one of our tasks will be 
to recover his estimate.  Suppose initially that $H$ 
is bounded from below and 
that $J$ is the lower part of its spectrum, 
$J = \{\lambda_1, \dots, 
\lambda_n\}$.  It is clear that if $f \ge 0$ 
is nondecreasing on $J$, then (4) 
produces an upper bound on $\Sigma f(\lambda_{j})$ when the first term 
on the right is dropped and the second sum is simplified by extending the 
sum to all $p$.  
Actually, a substantially stronger statement is obtainable:

\proclaim{Theorem 5}  
(asymptotically \# for Dirichlet Laplacian or power-law potentials)  Suppose 
that the same assumptions as in \rom{Theorem 1} apply.  
Let $f$ be a nonnegative 
function on $J = \{\lambda_1, ..., \lambda_n\}$ such that for some 
$z \in 
(\lambda_n,\lambda_{n+1}]$,  $f(\lambda)(z- \lambda)^{-2}$ is 
nondecreasing in $\lambda$ for $\lambda \in J$.  Then
$${d \over 4} \sum\limits_{j=1}^{n} f(\lambda 
_j)\le\sum\limits_{j = 1}^{n} {f(\lambda_j) \over 
z -\lambda_j}T_j.
\tag11$$
In particular, with $f(\lambda) = (z-\lambda)^2$, this implies that
$$
{d \over 4} \sum\limits_{j =1}^{n} (z -\lambda 
_j)^{2}\le \sum\limits_{j= 1}^{n} 
(z -\lambda_j)T_j. \tag 12$$
\endproclaim

\remark{Remarks} 
For the Dirichlet Laplacian with $z = \lambda_{n+1}$, 
(12) becomes a key 
formula of Yang's article ([Ya91], above his eq.\ (24) 
in the 1991 numbering 
or above (25) in the 1995 numbering), from which the subsequent bounds in that 
article follow readily.  As observed by Yang, 
when $n \rightarrow \infty$ 
 with $p=2$, both sides have the same 
 leading growth rate when we substitute 
the Weyl asymptotics for $\lambda_j$ in the case of the Dirichlet 
Laplacian.  On the other hand, if $T_j$ is replaced by $\sigma 
\lambda_j$, then inequality (12) may be interpreted as an inequality on 
a polynomial of the form 
$$
P_n(z) := nz^2- 2\left(1+{2\sigma  \over 
d}\right)\left(\sum\limits_{k=1}^{n} \lambda_k\right) 
z + {d+4\sigma  \over d}\sum\limits_{k = 1}^{n} \lambda_k^2.$$
This allows sharper estimates as follows:
\endremark

\proclaim{
Proposition 6} (\# for harmonic oscillator, $d=1$)
Let $H$ be as in \rom{Theorem 1} and suppose that the bound 
\rom{(10)} holds for some  
$\sigma\in (0,1]$.  
(Recall that $\sigma = \beta/(2+\beta)$ results from the 
virial assumption \rom{(9)} and that $\sigma=1$ 
applies to the Dirichlet Laplacian 
or a purely magnetic Schr\"odinger operator.)  Then 
\roster
\item"(i)" $P_n(z)  \le  0$ for all  $z \in [\lambda_n, \lambda_{n+1}]$.
\item"(ii)"
$D_n:= \left(\left(1+{2\sigma  \over d}\right){1 \over 
n}\sum\limits_{k=1}^{n} \lambda_k\right)^{2} - 
\left(1+{4\sigma  \over d}\right){1 \over n}\sum\limits_{k = 1}^{n} 
\lambda_k^2 \ge 0.$
\item"(iii)"$[\lambda_n, \lambda_{n+1}]\subseteq 
[\lambda_-(n),\lambda_+(n)], $
where 
$$\lambda_\pm(n) := \left(1+{2\sigma  \over d}\right) 
{1 \over n}\sum\limits_{k = 1}^{n} \lambda_k\pm  
\sqrt {D_n}.$$
In particular, for the gap $\lambda_{n+1}  - \lambda_{n}$, 
$$
\lambda_{n+1} - \lambda_n\le  2\sqrt {D_n}.$$
\endroster
\endproclaim

\demo{Proof of Proposition \rom{6}}
Statement (i) is a restatement of inequality (12).  
This implies that $n D_n = - 
\min(P_n(z))$ is nonnegative, giving (ii).  
Statement (iii) is merely the explicit 
determination of the upper and lower bounds.
\hfill\qed
\enddemo

\demo{Discussion}   
Observe that inequality (ii) goes in the opposite direction 
from the Cauchy-Schwarz inequality.  The underlying reason that it implies 
a simple but tight bound (iii) on the gap is that our technique is 
nonvariational, unlike most other gap estimates. 

For the Dirichlet Laplacian, $\lambda_j \sim j^{2/d}$ in accordance with 
Weyl's asymptotic formula, and this saturates our bound at leading order.  
Inequality (ii) is likewise asymptotically sharp for Schr\"odinger operators, 
in the sense that for Hamiltonians on $L^2(R^d)$ with power-law potentials,
$$H = |{\bold p}|^2 + |{\bold x}|^{\beta},$$
$\beta> 0$, such as the harmonic oscillator ($\beta=2$), 
the corresponding Weyl 
formula is
$$ \lambda_n \sim n^{2 \sigma /d} $$
[QuRo77], which again saturates the inequality for $D_n$ at leading order. 

Even more strikingly, we can compare the gap estimate from (iii) with the 
exact situation for the harmonic oscillator ($\beta=2$) in one dimension as 
follows.  The eigenvalues are $\lambda_n  = 2n-1$, and we calculate:
$$
\sigma = 1/2$$
$$
{1 \over n}\sum\limits_{k = 1}^{n} \lambda_k=
n $$
$$
{1 \over n}\sum\limits_{k = 1}^{n} \lambda_k^{2} =
 {1 \over 3}(4n^2 -1)$$
$$D_n= 1$$
and therefore $\lambda_n$ and $\lambda_{n+1}$ are exactly equal to the 
bounds 
$$
\lambda_\pm(n) = 2n\pm1. $$
\enddemo

\demo{Proof of Theorem \rom{5}}
We majorize the last term in (4) by:
$$\eqalign{
4\sum\limits_{{\scriptstyle\alpha \atop
\scriptstyle j\le n}\atop
\scriptstyle q\ge n+1} {f(\lambda_j) \over \lambda 
_q- \lambda_j}T_{\alpha jq}&\le 
4\sum\limits_{\alpha =1}^{d} \sum\limits_{j=1}^{n} 
{f(\lambda_j) \over z -\lambda_j} 
\sum\limits_{q = 1}^{\infty} T_{\alpha jq}- 
4\sum\limits_{\alpha =1}^{d} \sum\limits_{j\le n} 
{f(\lambda _j) \over z -\lambda_j}
\sum\limits_{q = 1}^{n} T_{\alpha jq}\cr
&=4
\sum\limits_{j = 1}^{n} {f(\lambda_j) 
\over z -\lambda_j}\,T_j - 4\sum\limits_{\alpha 
=1}^{d} \sum\limits_{j = 1}^{n} {f(\lambda_j) \over 
z - \lambda_j} \sum\limits_{q= 1}^{
n} T_{\alpha jq}  }$$
Comparing (11) with (4), it suffices to show that 
$$\sum\limits_{\alpha = 1}^{d} \left(-2\sum\limits_{
j,m=1}^{n} {f(\lambda_j)-f(\lambda 
_m) \over \lambda_j -\lambda_m}\,T_{\alpha 
jm} - 4\sum\limits_{j\le n}{f(\lambda 
_j) \over z - \lambda_j}\sum\limits_{m= 
1}^{n} T_{\alpha jm}\right) $$
is negative.  By symmetrizing the last term in this 
expression with respect to 
$j,m$, we get
$$\sum\limits_{\alpha  = 1}^{d} \sum\limits_{j ,m = 
1}^{n} \left(-2\,{f(\lambda_j) - 
f(\lambda_m) \over \lambda_j -\lambda_m}-
2\, {f(\lambda_j) \over z -\lambda_j} -2 
\,{f(\lambda_m) \over z - \lambda_m}\right)\,
T_{\alpha jm},$$
and some straightforward algebra shows that the expression in parentheses 
can be rewritten
$$\left(-2\left(z -\lambda_m\right)\left(z -\lambda_j\right)
{f(\lambda_j)(z -\lambda_j 
)^{-2} - f(\lambda_m)(z -\lambda 
_m)^{-2} \over (\lambda_j- \lambda 
_m)}\right),$$
which is negative by hypothesis.
\hfill\qed
\enddemo

A variant of these arguments, corresponding to a choice $J = 
\{\lambda_{n}, \dots, \lambda_{N}\}$ rather than $\{\lambda_{1}, \dots, 
\lambda_{N}\}$, allows $\lambda_{N+1}$ to be controlled by the 
eigenvalues in $J$ alone.  While the bounds do not appear as tight as in 
Propositions 6 and 7, they use less information:

\proclaim{Proposition 7}
  With the same assumptions as in \rom{Proposition 6}, but this 
time with
$$\eqalign{
P_{N,n}(z) :=P_N(z) -P_{N,n}(z) =(N-n) 
z^2&- 2\left(1+{2\sigma  \over d}\right)
\left(\sum\limits_{
k=n+1}^{N} \lambda_k\right) z\cr 
& +{d+4\sigma  \over 
d}\sum\limits_{k = n+1}^{N} \lambda_k^2
.}$$
\roster
\item"(i)"
$P_{N,n}(z)  \le  0$ for all  $z \in [\lambda_{N}, \lambda_{N+1}]$.
\item"(ii)" 
$D_{N,n} := \left(\left(1+{2\sigma  \over d}\right) {1 
\over N-n} \sum\limits_{k=n+1}^{N} \lambda_k\right) 
^2-\left(1+{4\sigma  \over d}\right) {1 \over N-n} 
\sum\limits_{k = n+1}^N\lambda_k^2\ge  0.$
\item"(iii)"
$[\lambda_N,\lambda_{N+1}]\subseteq  
[\lambda_-(N,n),\lambda_+(N,n)],$
where 
$$\lambda_\pm (N,n) :=\left(1+{2\sigma  \over d}\right) 
{1  \over N-n} \sum\limits_{k = n+1}^N \lambda_k
\pm  \sqrt {D_{N,n}}. $$
In particular, for the gap $\lambda_{N+1}  - \lambda_{N}$, 
$$
\lambda_{N+1} -\lambda_N\le 2\sqrt {D_{N,n}}.$$
\endroster
\endproclaim

\demo{Proof}  
In identity (7) choose 
$J = \{\lambda _{n+1}, \dots, \lambda _{N}\}$ and 
substitute the formula (9) in the first term on the right.  
The analogue of 
(12),
$$
{d \over 4}\sum\limits_{j = n+1}^{N} (z- 
\lambda_j)^2\le \sum\limits_{j = 
n+1}^{N} (z - \lambda_j)T_j,$$
follows because the final term in (7), 
$$4\sum\limits_{\scriptstyle\lambda_j\in J\atop
\scriptstyle\lambda_q\in J^c} \left({(z -\lambda 
_j)(z-\lambda_q) \over \lambda_q- 
\lambda_j}\right)\sum\limits_{\alpha }T_{\alpha 
jq}\le 0 $$
for $z \in [\lambda_{N}, \lambda_{N+1}]$.   
(In this sum, either $q > N  \ge  
j$ or $q  \le  n < j$).
\hfill\qed\enddemo

\head III.  Bounds for moments and partition functions\endhead

In order to generate further useful inequalities from the 
identity (4), we shall 
assume the existence of a function $r(x)$ such that for $
x,y \in J$, 
$${f(x) - f(y) \over x-
y} \ge  {r(x) + r(y) \over 2}\leqno \text{(H1)} $$
The existence of such a function $r$ 
is of course a restriction on $f(x)$, which in 
practice is that $f'''  \le  0$:


\proclaim{Lemma 8}  
Hypothesis \rom{(H1)} holds with $r(s) = f '(s)$ provided that 
$f \in C^3(I)$ 
for some interval $I$ containing $J$ and that 
$f '''(s)  \le  0$ for all $s \in I$.
\endproclaim

\demo{Proof}  
Without loss of generality we can rescale $f$ and subtract a constant 
to set $y = 0$, 
$x = 1$, and $f(0) = 0$, with $f '''(s)  \le  0$ 
for all $s \in [0,1]$.  The derivative $f '$ is 
concave on that interval, and therefore 
$f '(s)  \ge  s f '(0) + (1-s) f '(1)$.  If 
we integrate this inequality from 0 to 1, we obtain
$$
f(1) \ge  {1 \over 2}\, f^\prime(0) + {1 \over 2}\, f^\prime(1)
$$
as required.
\hfill\qed
\enddemo

\remark{Remarks} %%Remarks.  
In what follows we shall generally choose $r(x) = f '(x)$, since 
this seems to produce the sharpest corollaries.  Complementary bounds in 
the other direction result in principle from the alternative hypothesis 
$$
 {f(x) - f(y) \over x-
y} \le  {R\left({x}\right) + R\left({y}\right) \over 2}\leqno \text{(H2)} $$
(just replace $r$ by $R$ and reverse the inequality in (13), below), 
but we have 
not found them useful.  If $f$ is a quadratic or linear polynomial, both 
hypotheses hold with equality for $r(x) = R(x) = f '(x)$ 
(cf.\ Proposition 2).
\endremark

\proclaim{Theorem 9}
 With assumptions as in \rom{Theorem 1} and inequality \rom{(H1)}, 
$${d \over 2}\sum\limits_{\lambda_j\in J} 
f(\lambda_j) \le -\sum\limits_{\lambda_j\in J}  
r(\lambda_h)T_j+ 
\sum\limits_{\scriptstyle\lambda_j\in J\atop
\scriptstyle\lambda_q\in J^c}  \left({2f(\lambda 
_j) \over \lambda_q-\lambda_j}+  
r(\lambda_j)\right)\sum\limits_\alpha T_{\alpha jq}  .\tag 13$$                       
In particular, for $f(\lambda ) = (z-\lambda )^p$ with 
$p  \ge  2$ and $z \in [\lambda_{n}, 
\lambda_{n} + p(\lambda_{n+1}- \lambda_{n})/2]$,
$${d \over 2p}\sum\limits_{j = 1}^{n} (z - \lambda 
_j)^p \le  \sum\limits_{j=1}^{n} 
(z - \lambda_j)^{p-1}T_j . 
\tag 14    
$$
\endproclaim

\demo{Proof}  
Apply hypothesis (H1) to identity (4), recalling that $T_{\alpha j q}  
\ge  0$. The result is
$$\sum_{\lambda_j\in J}f(\lambda_j)\le -2 
\sum_{\scriptstyle
\lambda_j,\lambda_m\in J\atop\scriptstyle \lambda_j\ne \lambda_j}
{r(\lambda_j)+ r(\lambda_m)\over 2}\,T_{\alpha jm} +4
\sum_{\scriptstyle\lambda_j\in J\atop\scriptstyle \lambda_q\in J^c}
{f(\lambda_j)\over \lambda_q-\lambda_j}  \, T_{\alpha jq}.$$
Because the first sum on the right is symmetric in 
$j,m$, we may rewrite it as:
$$
\sum\limits_{\scriptstyle\lambda_j,\lambda_m\in J\atop
\scriptstyle\lambda_j\ne\lambda_m}  {r(\lambda_j 
) +r\left(\lambda_m\right) \over 2}\,T_{\alpha 
jm} =  \sum\limits_{\scriptstyle\lambda_j,\lambda 
_m\in J\atop
\scriptstyle\lambda_j\ne \lambda_m}  r\left(\lambda_j 
\right)T_{\alpha jm}.$$
Since 
$$-2\sum\limits_{\scriptstyle\lambda_j,\lambda_m\in 
J\atop
\scriptstyle\lambda_j\ne \lambda_m}  r\left(\lambda_j 
\right)T_{\alpha jm} = -2 \sum\limits_{\scriptstyle\lambda 
_j\in J\atop
\hbox{all }\scriptstyle\lambda_m}  r\left(\lambda_j 
\right)T_{\alpha jm} +2 \sum\limits_{\scriptstyle\lambda_
j\in J\atop
\scriptstyle\lambda_m\in J^c}  r\left(\lambda_j 
\right)T_{\alpha jm}   $$
and 
$$\sum\limits_{\scriptstyle\alpha \atop
\hbox{all }\scriptstyle\lambda_m}  T_{\alpha jm} = 
T_j,   $$
by (5), inequality (13) follows after summing on $\alpha$ and dividing by 
2. 

For the moment inequality (14), notice that $f(\lambda) := (z - \lambda)^p$ 
satisfies (H1) for all $\lambda < z$  by Lemma 8.  The final term in (13) is 
negative for the given range of $z$ 
because of the positivity of $T_{\alpha j q} $
and the observation that for $q  \ge  n+1$, its coefficient is
$$\eqalign{
\Bigl(2\left(z - \lambda_j\right)&- p\left(\lambda_q- 
\lambda_j\right)\Bigr)\left(z - \lambda_j\right)^{p-1} \cr
&\le  \left(2\left(z - \lambda_j\right) - p\left(\lambda 
_{n+1} - \lambda_j\right)\right)\left(z - \lambda_j 
\right)^{p-1} \le  0.}\eqno{\qed}$$
\enddemo

One of the most significant functions in physics is the partition 
function, or 
trace of the heat operator,
$$        Z(t) := \text{tr}(\exp(-tH)),                                                 
$$
and evidently our methods apply with $f(\lambda) = \exp(-\lambda t)$, $t 
> 0$,  and  $J =$ 
the entire spectrum.  The set $J$ may be allowed to be the entire 
spectrum assuming that the eigenvalues grow sufficiently fast, which is 
guaranteed by the usual Weyl asymptotics for Dirichlet Laplacians or 
Schr\"odinger operators on finite domains, or for Schr\"odinger operators 
on infinite domains if the potential grows sufficiently at infinity, for 
example.  We henceforth assume the convergence of $Z(t)$.  

Since $f(\lambda)$ has a negative third derivative in its variable 
$\lambda$ for any $t>0$, hypothesis (H1) is satisfied, and hence from (13), 
$$
Z\left(t\right) \le  -\left({2 \over d}\right) \sum\limits_{
j}  \left(-t \exp \left(-t \lambda_j\right)\right) 
T_j = \left({2 t \over d}\right) \sum\limits_{j}  
{T_j \over {\lambda }_{j}} \left(\exp \left(-t \lambda 
_j\right)\right) \lambda_j    \tag 15  $$
As we have seen, in many cases the factor $T_{j} /\lambda_{j}$ can 
be replaced by a constant $\sigma$.  Thus for the Dirichlet Laplacian on a 
bounded domain, which corresponds physically to the ideal gas and 
$\sigma=1$, (14) yields 
$$Z_{i.g.}(t)\le -{2t\overwithdelims() d}Z'_{i.g.}(t)\tag 16
$$
which implies:

\proclaim{Corollary 10} 
(\# for ideal gas = Dirichlet Laplacian, for large volume or 
$t\rightarrow 0$)  $Z(t) t^{d/2}$ is a nonincreasing function.
\endproclaim

\demo{Proof}
We can rewrite (16) as 
$$                              
{d\over dt}\,\ln (Z_{i.g.}(t))\le -{d\overwithdelims () 2}
{d\ln (t)\over dt},$$
or, equivalently,
$${d\over dt}\,\ln (Z_{i.g.}(t)t^{d/2})\le 0.\tag 17$$
\hfill\qed\enddemo


This should be compared with the known result from comparison of Green 
functions (see, e.g., [Da]), that 
$$
Z_{i.g.}(t)  \le  {\text{Vol}(\Omega)\over
(4 \pi t)^{d/2}},
$$
which also follows from integrating (16) between 0 and $t$, 
using the fact that 
$$\lim_{t\to 0} Z_{i.g.}(t)(4\pi t)^{d/2} =\text{vol}(\Omega).$$
With the virial theorem, if (9) holds, we obtain:

\proclaim{Corollary 11}  
(\# as $t \rightarrow 0$)  Let $H =   |{\bold p}|2 + 
V({\bold x})$ on $L^2(R^d)$, 
where ${\bold x}\cdot \nabla V   \le  
\beta V$, and suppose that the spectrum of $H$ is 
discrete and $Z(t)$ convergent for $t>0$.  Then 
$$
Z(t)\le -{2\beta t\overwithdelims () (2+\beta)d} Z'(t),$$
so 
$$
Z(t)t^{(2+\beta)d/2\beta}\text{ is a nonincreasing function.}
$$
\endproclaim

\demo{Proof}  
Substitute for $T_j/\lambda_{j}$ as in Corollary 4  with (10) and 
argue as for Corollary 10.     
\hfill\qed\enddemo 

\proclaim{Corollary 12}  
(\# for harmonic oscillator)\/ In addition to the assumptions 
of \rom{Corollary 11}, suppose that there is a minimal spectral gap 
$\omega  \le  
|\lambda_{k} - \lambda_{j}|$ for all $\lambda_{j} \ne
\lambda_{k}$  (exact degeneracies are allowed).  Then
$${Z} \left({t}\right) \le  - \left({4\beta  \tanh \left({\omega  
\over 2}t\right) \over \omega d\left(2+\beta \right)}\right) 
{Z}^{\prime}\left({t}\right),          \tag 18   $$
and 
$
{Z} \left({t}\right) {\left({\sinh \left({{\omega  \over 
2}t}\right)}\right)}^{(2+\beta )d/2\beta }    $
is a nonincreasing function.
\endproclaim

\demo{Proof}
For the partition function, with 
$J = $the entire spectrum, identity (4) becomes
$$\eqalign{
\sum\limits_k  \exp \left(- t \lambda _k\right) &= 
-2\sum\limits_{j \ne  k}  {\exp \left(- t \lambda_j 
\right) - \exp \left(- t \lambda_k\right)  \over \lambda_j 
- \lambda_k}\,T_{\alpha jk} \cr
&= \sum\limits_{j \ne  k}  (\exp (- t \lambda 
_j)+ \exp (- t \lambda_k))\times\cr
&\qquad  \left(
{\exp \left(- t \lambda_j\right) - \exp \left(- t \lambda 
_k\right)  \over \left({\lambda_k - \lambda_j \over 
2}\right)\left(\exp \left(- t\lambda_j\right) + \exp \left(- t 
\lambda_k\right)\right)}\right)T_{\alpha jk}\cr
&= t\sum\limits_{j \ne  k}  2 \exp \left(- t \lambda 
_j\right){\tanh \left({\lambda_k - \lambda_j \over 
2}\,t\right) \over \left({\lambda_k - \lambda_j \over 
2}\,t\right)}\,T_{\alpha jk} \cr
&\le t\sum_{j\ne k} 2\exp(-t\lambda_j){\tanh \left({\omega\over 2} 
\,t\right)\over \left({\omega\over 2}\,t\right)}\, T_{\alpha jk},}
$$
using the symmetry in $j,k$ and the noting that 
$\tanh(x) /x$ is a decreasing 
function.  When  we sum over $\alpha$, we get
$$
d Z(t) \le  \sum\limits_{j ,k,\alpha }  2 \exp \left({- t 
{\lambda }_{j}}\right){\tanh \left({ \omega  \over 
2}\,t\right) \over \left({{\omega  \over 2}}\right)}{T}_{\alpha jk} 
= \sum\limits_{j}  2 \exp \left({- t {\lambda 
}_{j}}\right){\tanh \left({ {\omega  \over 2}t}\right) \over 
\left({{\omega  \over 2}}\right)}{T}_{j}    $$
so, with (10), 
$$d {Z} \left({t}\right) \le  -{4 \over \omega } \tanh 
\left({{\omega  \over 2}t}\right) {Z}^{\prime}(t) {\beta  \over 
2+\beta }$$
This is equivalent to (18), which in turn may be written
$$0\ge {d\over dt}\,\left( Z(t) \left(\sinh \left({\omega \over 2}\,
r\right)\right)^{d(2+\beta)/2\beta}\right).\eqno{\qed}$$
\enddemo

Corollary 12 is again optimal, for the exact partition function for the 
harmonic oscillator, normalized so that 
$\omega=2$, $\beta=2$ ($\sigma=1/2$), 
is 
$${Z}_{h.o.}\left({t}\right)  = 
{\left({2 \sinh \left({t}\right)}\right)}^{-
d}.$$
Inequalities similar to these corollaries hold for the 
finite partition function
$${Z}_{N}\left({t}\right) := \sum\limits_{j = 1}^{N} \exp 
\left({-t{\lambda }_{j }}\right) $$
provided that $t$ is sufficiently large (viz., $>2/(\lambda_{N+1}  - 
\lambda_{N})$).  In addition, some slightly more refined estimates for the 
derivatives of $Z$ are obtainable, since for any function 
$f(\lambda)$ with 
$f '''  \le  0$ and $J =$ the entire spectrum, inequality (13) becomes
$${d \over 2}\sum\limits_{k}  f\left({{\lambda }_{k}}\right) \le 
 -\sum\limits_{k}  f'\left({{\lambda }_{
k}}\right){T}_{k},    
\tag 20$$
assuming convergence. For example, choosing 
$$
	   f(\lambda ;t) = 3+\lambda t \exp(- \lambda  t)  
$$
(the 3 ensures that $f '''  \le  0$ for all $t$), 
leads to a second order differential 
inequality,
$$
{t}^{2} Z''(t) + \left({{d \over 2}-2}\right) t Z' (t) - 
{3d \over 2} Z(t) \ge  0.  $$

{\bf Acknowledgments.}  
We are grateful to P.~Kr\"oger for bringing the 
article by Yang to our attention and to him and to J.~Fleckinger for helpful 
conversations.  We also wish to acknowledge the hospitality of the Erwin 
Schr\"odinger Institute in Vienna, where our collaboration began.

\Refs

\widestnumber\key{PAPOWE56}
\ref\key AsBe  
\by Ashbaugh, Mark S., Benguria, Rafael D.\paper A sharp bound for the 
ratio of the first two eigenvalues of Dirichlet Laplacians and extensions 
\jour Ann. Math. \vol 135\yr 1992\pages 601--628\endref
\ref\key Da 
\by Davies, E.B. \paper Heat Kernels and Spectral Theory
\inbook Cambridge Tracts in 
Mathematics \vol 92 \publaddr Cambridge  \publ Cambridge University Press 
\yr 1989\endref
\ref\key EaKa  
\by Eastham, M.S.P., Kalf, H.\paper Schr\"odinger-Type Operators with 
Continuous Spectra \inbook Research Notes in Mathematics \vol 65 
\publaddr Boston \publ Pitman
\yr 1982\endref
\ref\key Ha88   
\by  Harrell, Evans M., II\paper General Bounds for the Eigenvalues of 
Schr\"odinger Operators \pages 146--166 \inbook Maximum Principles and 
Eigenvalue Problems in Partial Differential Equations 
\bookinfo P.W. Schaefer, 
editor  \publaddr Essex, England  
\publ Longman House \finalinfo and New York, Wiley, 1988\endref
\ref\key Ha93   
\by Harrell, Evans M., II \paper Some geometric bounds on eigenvalue gaps 
\jour Communications in Partial Differential Equations 
\vol 18\yr 1993\pages 179--198\endref
\ref\key HaMi95   
\by Harrell, Evans M., II and Michel, Patricia L.\paper Commutator 
Bounds for Eigenvalues, with Applications to Spectral Geometry 
\jour Communications in Partial Differential Equations
\vol 19\yr 1994\pages 2037--2055
\moreref Erratum \vol 20\yr 1995\pages 1453\endref
\ref\key HiPr80  
\by Hile, G.N., M. H. Protter\paper Inequalities for eigenvalues of the 
Laplacian\jour Indiana Univ. Math. J. \vol 29 \yr 1980
\pages 523--538\endref
\ref\key Ho90  
\by Hook, Stephen M.  \paper Domain-independent upper bounds for 
eigenvalues of elliptic operators\jour Trans. Amer. Math. Soc. 
\vol 318 \yr 1990 
\pages 615--642\endref
\ref\key PaPoWe56  
\by Payne, L.E., P\'olya, G., Weinberger, H.F.\paper On the ratio of 
consecutive eigenvalues\jour J. Math. and Phys. \vol 35 
\yr 1956\pages 289--298\endref
\ref\key QuRo77\by Quigg, C.,
Rossner, J. L.\paper Quantum Mechanics with applications to quarkonia
\jour Physics Reports \vol 56 \yr 1977\pages 167--235\endref
\ref\key Ya91  
\by Yang, Hongcang\paper An estimate of the difference between consecutive 
eigenvalues \paperinfo
preprint IC/91/60 of the Intl. Centre for Theoretical Physics, 
\yr 1991  \finalinfo Revised preprint, from Academia Sinica, 1995\endref
\endRefs

\enddocument

